% FILE: 07_open_questions_and_next_work.tex
% PROJECT: otel-weaver-exp-001-custom-pi-registry
% THESIS ROLE: otel-semconv-evidence-surface-experiment

\section{Open Questions and Next Work}
\label{sec:otel-weaver-exp-001-custom-pi-registry-questions}

\subsection{Known Gaps}
\label{sec:otel-weaver-exp-001-custom-pi-registry-questions-gaps}

The following gaps are identified from the evidence audit and marked by epistemic status:

\textbf{GAP-OW-001: Receipt not materialized.}
The BLAKE3 receipt declared in \texttt{otel-weaver-source.manifest.yaml} points to
\texttt{/Users/sac/process-intelligence/receipts/otel-weaver-source-receipt.json}. This file
has not been confirmed to exist on disk. A declared receipt and a materialized receipt are
structurally distinct: the former is an aspiration; the latter is evidence. The PARTIAL status
of EXP-001 is primarily attributable to this gap.

\textbf{GAP-OW-002: Registry directory not confirmed.}
The \texttt{weaver registry validate} command requires a \texttt{registry/} directory. The
\texttt{otel-weaver/registry/} path does not appear to exist in the confirmed evidence.
Without this directory, the Weaver CLI cannot validate \texttt{process\_pi.yaml}, and the
machine-verifiability claim is unconfirmed (FUTURE WORK).

\textbf{GAP-OW-003: Test scope ambiguity.}
The integration test \texttt{test\_weaver\_integration\_synthesis} in
\texttt{sources/wasm4pm/tests/weaver\_integration\_tests.rs} may validate only the Rust
bridge/validator logic independently of \texttt{process\_pi.yaml}. If the test does not
exercise the schema attributes directly --- by constructing spans with required fields
present or absent and verifying the corresponding \texttt{Admission}/\texttt{Refusal} outcome
--- then the test suite does not provide evidence for EXP-001's feedstock schema claim.

\textbf{GAP-OW-004: \texttt{process.pi.activity.type} not enforced.}
The \texttt{process.pi.activity.type} attribute is recommended, not required. The current
\texttt{validator.rs} \texttt{RefusalCode} enum does not include a
\texttt{MissingActivityType} variant, which means a span without an architectural category
tag can pass the admission gate and reach the conformance court without being classified
by activity type. Whether this is a design choice or an oversight is not determinable from
the evidence corpus alone (INTERPRETATION).

\textbf{GAP-OW-005: No pm4py-layer validation.}
The ALIVE checkpoint was self-certified by the ggen manufacturing authority. No independent
pm4py-based fitness or precision check has been run against event logs produced from
\texttt{process.pi.activity}-instrumented spans. Schema compliance at the feedstock layer
has not been demonstrated to propagate to conformance quality at the replay layer.

\subsection{Deferred Gates}
\label{sec:otel-weaver-exp-001-custom-pi-registry-questions-deferred}

The following gates are declared in the architecture but deferred pending resolution of the
known gaps:

\begin{itemize}
  \item \textbf{Weaver CLI validation gate} --- requires materialization of the
        \texttt{registry/} directory and confirmed execution of \texttt{weaver registry validate}
        without errors. Blocks machine-verifiability claim for \texttt{process\_pi.yaml}.
  \item \textbf{Receipt materialization gate} --- requires the BLAKE3 receipt file to exist at
        the declared store location and for its contents to be verifiable against the
        \texttt{process\_pi.yaml} artifact. Blocks full ALIVE status.
  \item \textbf{Schema-to-replay propagation gate} --- requires a pm4py conformance check
        against an event log produced from live spans. Blocks the claim that feedstock schema
        compliance propagates to process-mining-layer conformance quality.
\end{itemize}

\subsection{Research Risks}
\label{sec:otel-weaver-exp-001-custom-pi-registry-questions-risks}

\textbf{Risk 1: Schema version lock.}
The registry is pinned to \texttt{schema\_url: https://opentelemetry.io/schemas/1.25.0}. If the
OTel standard evolves in ways that deprecate or rename required attribute types, the registry
will require update, and the BLAKE3 receipt chain will need to be re-sealed. The \texttt{BridgeRx}
invariant (no residual from bridging) may not hold across breaking OTel schema changes.

\textbf{Risk 2: Recommended-field gap exploitation.}
The recommended attributes (\texttt{process.pi.token.state\_before},
\texttt{process.pi.token.state\_after}, \texttt{process.pi.activity.type}) are not enforced
by the admission gate. A span that omits all three can pass the gate and reach the conformance
court. For Petri net replay, a span without token state information provides weaker evidence
than one with it. The court may still replay the span, but with lower precision and possibly
inflated fitness. The degree of degradation has not been quantified (FUTURE WORK).

\textbf{Risk 3: Self-certification circularity.}
All three checkpoints were issued by the ggen manufacturing authority. The authority that
manufactured the registry is the same authority that certified its ALIVE status. An independent
verifier --- a human auditor, a CI pipeline, or a pm4py replay --- has not confirmed any of the
gateway assertions. This circularity is the primary research risk for the dissertation's
evidence-chain integrity claim as it applies to EXP-001.

\subsection{Next Receipted Work}
\label{sec:otel-weaver-exp-001-custom-pi-registry-questions-next}

The following work items, if completed and receipted, would advance EXP-001 from PARTIAL to
ALIVE status:

\begin{enumerate}
  \item Materialize the \texttt{registry/} directory from \texttt{process\_pi.yaml} and run
        \texttt{weaver registry resolve} and \texttt{weaver registry validate}. Record the
        output in a new checkpoint file. Seal the checkpoint with a BLAKE3 receipt at
        \texttt{receipts/otel-weaver-source-receipt.json}.
  \item Extend \texttt{test\_weaver\_integration\_synthesis} or add a new test that constructs
        spans with each required attribute absent in turn, verifies the corresponding
        \texttt{RefusalCode} is returned, and verifies that spans with all required attributes
        present produce a valid \texttt{Admission} value. Commit the test with a
        \texttt{research-pm4py} scope commit and record the test run receipt.
  \item Run a pm4py conformance check: instrument a minimal process (at minimum: a two-activity
        sequence with a scheduled and a completed lifecycle event) using spans that satisfy the
        \texttt{process.pi.activity} schema, project the spans to OCEL using
        \texttt{OtelSpanToOcelEventProjection}, and measure fitness and precision against the
        declared Petri net model. Record the pm4py output in a \texttt{checkpoint} file and
        seal it. A fitness $> 0.8$ and precision $> 0.8$ against a simple reference model
        would constitute sufficient evidence for the propagation claim.
  \item Add a \texttt{MissingActivityType} variant to the \texttt{RefusalCode} enum in
        \texttt{validator.rs} and make \texttt{process.pi.activity.type} required rather than
        recommended, if the process intelligence court requires architectural category
        classification for all admitted spans. Document the decision in a doctrine addendum
        (FUTURE WORK / AUTHOR THESIS).
\end{enumerate}
